% GENERATED FILE. Edit canonical lessons, then run npm run generate:books.
% canonical-source-hash: fnv1a64:2e01c91836f60fab
% canonical-lessons: MR-C176-dhagal, MR-C176-pavsali, MR-C176-sakya, MR-C176-asakya, MR-C176-mukhya

\chapter{Cloudy, Rainy, Possible, Impossible, Main}
\label{ch:mr-cloudy-rainy-possible-impossible-main}
\hlchaptermodality{176}

\begin{chapteropening}
\textbf{By the end of this chapter:} \emph{I can say cloudy, rainy, possible, impossible and main.}

Complete the last lesson of chapter 176: I can say cloudy, rainy, possible, impossible and main.
\end{chapteropening}

\section[\texorpdfstring{\mr{ढगाळ} (ḍhagāḷ)}{ढगाळ (ḍhagāḷ)}]{\mr{ढगाळ} (ḍhagāḷ) — cloudy}

\label{lesson:MR-C176-dhagal}



\begin{quote}
Before the new one: say the Marathi for proper, then the Marathi for wrong.
\end{quote}

\subsection*{You'll want to know: \mr{ढगाळ}}
\textbf{\mr{ढगाळ}} — \emph{ḍhagāḷ} — \textquotedblleft{}cloudy\textquotedblright{}.

\begin{grammarlens}[title={What you've built}]
One more word for this chapter.
\end{grammarlens}

\subsection*{Your turn}
\begin{itemize}
\raggedright
  \item \emph{Say it:} \emph{ḍhagāḷ}
  \item \emph{Say it:} \emph{ḍhagāḷ}, once more
  \item \emph{Say it:} say \emph{ḍhagāḷ}
  \item \emph{Recall:} say the Marathi for proper, then the Marathi for wrong, then say \emph{ḍhagāḷ} again
\end{itemize}

\begin{tcolorbox}[breakable,colback=teal!4,colframe=teal!35!black,title={Before you move on}]
What does \mr{ढगाळ} mean? (\textquotedblleft{}Cloudy\textquotedblright{}.) Say it once more.
\end{tcolorbox}
\section[\texorpdfstring{\mr{पावसाळी} (pāvsāḷī)}{पावसाळी (pāvsāḷī)}]{\mr{पावसाळी} (pāvsāḷī) — rainy}

\label{lesson:MR-C176-pavsali}



\begin{quote}
Before the new one: say the Marathi for wrong, then the Marathi for cloudy.
\end{quote}

\subsection*{You'll want to know: \mr{पावसाळी}}
\textbf{\mr{पावसाळी}} — \emph{pāvsāḷī} — \textquotedblleft{}rainy\textquotedblright{}.

\begin{grammarlens}[title={What you've built}]
One more word for this chapter.
\end{grammarlens}

\subsection*{Your turn}
\begin{itemize}
\raggedright
  \item \emph{Say it:} \emph{pāvsāḷī}
  \item \emph{Say it:} \emph{pāvsāḷī}, once more
  \item \emph{Say it:} say \emph{pāvsāḷī}
  \item \emph{Recall:} say the Marathi for wrong, then the Marathi for cloudy, then say \emph{pāvsāḷī} again
\end{itemize}

\begin{tcolorbox}[breakable,colback=teal!4,colframe=teal!35!black,title={Before you move on}]
What does \mr{पावसाळी} mean? (\textquotedblleft{}Rainy\textquotedblright{}.) Say it once more.
\end{tcolorbox}
\section[\texorpdfstring{\mr{शक्य} (śakya)}{शक्य (śakya)}]{\mr{शक्य} (śakya) — possible}

\label{lesson:MR-C176-sakya}



\begin{quote}
Before the new one: say the Marathi for cloudy, then the Marathi for rainy.
\end{quote}

\subsection*{You'll want to know: \mr{शक्य}}
\textbf{\mr{शक्य}} — \emph{śakya} — \textquotedblleft{}possible\textquotedblright{}.

\begin{grammarlens}[title={What you've built}]
One more word for this chapter.
\end{grammarlens}

\subsection*{Your turn}
\begin{itemize}
\raggedright
  \item \emph{Say it:} \emph{śakya}
  \item \emph{Say it:} \emph{śakya}, once more
  \item \emph{Say it:} say \emph{śakya}
  \item \emph{Recall:} say the Marathi for cloudy, then the Marathi for rainy, then say \emph{śakya} again
\end{itemize}

\begin{tcolorbox}[breakable,colback=teal!4,colframe=teal!35!black,title={Before you move on}]
What does \mr{शक्य} mean? (\textquotedblleft{}Possible\textquotedblright{}.) Say it once more.
\end{tcolorbox}
\section[\texorpdfstring{\mr{अशक्य} (aśakya)}{अशक्य (aśakya)}]{\mr{अशक्य} (aśakya) — impossible}

\label{lesson:MR-C176-asakya}



\begin{quote}
Before the new one: say the Marathi for rainy, then the Marathi for possible.
\end{quote}

\subsection*{You'll want to know: \mr{अशक्य}}
\textbf{\mr{अशक्य}} — \emph{aśakya} — \textquotedblleft{}impossible\textquotedblright{}.

\begin{grammarlens}[title={What you've built}]
One more word for this chapter.
\end{grammarlens}

\subsection*{Your turn}
\begin{itemize}
\raggedright
  \item \emph{Say it:} \emph{aśakya}
  \item \emph{Say it:} \emph{aśakya}, once more
  \item \emph{Say it:} say \emph{aśakya}
  \item \emph{Recall:} say the Marathi for rainy, then the Marathi for possible, then say \emph{aśakya} again
\end{itemize}

\begin{tcolorbox}[breakable,colback=teal!4,colframe=teal!35!black,title={Before you move on}]
What does \mr{अशक्य} mean? (\textquotedblleft{}Impossible\textquotedblright{}.) Say it once more.
\end{tcolorbox}
\section[\texorpdfstring{\mr{मुख्य} (mukhya)}{मुख्य (mukhya)}]{\mr{मुख्य} (mukhya) — main, chief}

\label{lesson:MR-C176-mukhya}



\begin{quote}
Before the new one: say the Marathi for possible, then the Marathi for impossible.
\end{quote}

\subsection*{You'll want to know: \mr{मुख्य}}
\textbf{\mr{मुख्य}} — \emph{mukhya} — \textquotedblleft{}main, chief\textquotedblright{}.

\begin{grammarlens}[title={What you've built}]
One more word for this chapter.
\end{grammarlens}

\subsection*{Your turn}
\begin{itemize}
\raggedright
  \item \emph{Say it:} \emph{mukhya}
  \item \emph{Say it:} \emph{mukhya}, once more
  \item \emph{Say it:} say \emph{mukhya}
  \item \emph{Recall:} say the Marathi for possible, then the Marathi for impossible, then say \emph{mukhya} again
\end{itemize}

\begin{tcolorbox}[breakable,colback=teal!4,colframe=teal!35!black,title={Before you move on}]
What does \mr{मुख्य} mean? (\textquotedblleft{}Main, chief\textquotedblright{}.) Say it once more.
\end{tcolorbox}
